\chapter{Three Temporal Directions and Causal Reduction}
\label{ch:ch12_causal}

The question of this chapter is whether the Lorentzian four-plane used in every earlier chapter, and with it the light cone and the
distinction between past and future, can be obtained from the split geometry rather than postulated.

\paragraph{Motivation and scope.} This investigation originated in the intuition that observable causal order might follow from the
collective dynamics of the underlying system, rather than enter through an independently prescribed choice of a Lorentzian sector. The
quantum motivation comes from the sum-over-histories formulation, in which complex probability amplitudes associated with admissible
histories are combined before probabilities are extracted \cite{Feynman48}. Applied to the Spacetime--Electromagnetic Framework, this
suggests asking whether quantum dynamics and physical-state reduction associated with the split-signature geometry could determine an
effective Lorentzian causal structure.

This possibility is a research conjecture. Summing amplitudes does not by itself derive a spacetime signature, select a physical four-plane,
or establish microcausality. A complete realisation would require a specified space of histories, an action and integration prescription,
a consistent treatment of constraints, and a positive physical state space. In particular, the Lorentzian cone whose emergence is sought
cannot already be imposed through the admissibility conditions.

The present chapter investigates prerequisites and obstructions to this programme through explicit models. Its distinction between weak
and strong emergence (\cref{sec:baseline}) separates consistency within a selected Lorentzian sector from a dynamical explanation of that
selection. The conditional scalar reduction of \cref{cor:weak} achieves the former; the proposed emergence of causal order from the combined
quantum dynamics remains open, and the obstructions below hold for the models and assumptions in which they are proved.

The chapter starts from the free field on
$\R^{3,3}$, which for unrestricted data is ill-posed and has no time orientation (\cref{sec:baseline}); identifies the choice of a four-plane with the choice of a
complex structure and recovers one nappe of a light cone from twisted squares of spinors (\cref{sec:selector}); shows that no unconstrained
$SO(3,3)$-invariant adjoint scalar order parameter selects that choice cleanly (\cref{sec:typesel}); and derives the surplus-time constraints from a
gauge symmetry of the worldline that is chosen, not derived (\cref{sec:constraints}). \Cref{sec:tau} then analyses the clock reduction of
\cref{assum:clockred}, the two parameters called $\tau$ in this manuscript, and what the reduction keeps and removes; \cref{sec:universaltau} states
the universal-evolution hypothesis as an open hypothesis, with the ingredients a proposal would need. \Cref{sec:causalverdict} collects the
outcome. The chapter uses the clock and the interior four-plane of \cref{ch:ch02_lorentz} and the normal-plane model of \cref{ch:ch11_curvature};
the compatibility chapter (\cref{ch:ch13_compatibility}) uses its verdicts in the spacetime chain.

Several letters carry more than one meaning in this chapter, each fixed by its context: $B$ is the order parameter of \cref{sec:typesel}, the
Kaluza--Klein potentials $B^u,B^v$ of \cref{thm:ghostvec} and, in bold, the magnetic field; $C$ is the selector of \cref{sec:selector} and the
function $C=A_\mu h^{\mu\nu}k_\nu$ of \cref{thm:wlgauge}; $N$ is the normal plane and the lapse; $n$ is the null rotation of \cref{thm:flat} and a
lattice vector in \cref{ex:tachyons}; and $K=\tfrac12e_1e_4$ in \cref{thm:flat} is distinct from the grading root $\cK$ and from the kernels $K_{\psi\chi}$.
The letters $a,\dots,h$, $\kappa$, $\lambda$, $\lambda_2$ and $m$ in \cref{thm:degree} and its proof are coefficients local to that proof, and the
degree-eight function there is written $F$.

% --- source: tex/s27_baseline#0 ((head))


% --- source: tex/s27_baseline#1 (Causal structure: the six-dimensional baseline)
\section{Causal structure: the six-dimensional baseline}
\label{sec:baseline}

Every construction so far has used a Lorentzian four-plane chosen inside $\R^{3,3}$: the interior plane $E$ of the Dirac structure, or the slice $\Sigma$
of the curved model. Two levels must be distinguished. \emph{Weak emergence} means that an already selected $3{+}1$ sector carries a
consistent quantum causal structure. \emph{Strong emergence} means that the dynamics determine why that sector, with that causal structure, arises from the
six-dimensional theory. The light cone must not be put into the admissibility conditions and then be said to emerge. We fix the failure conditions in advance:
uncontrolled growing modes in the six-dimensional theory; a projector onto the four-dimensional sector imposed by hand; ghosts or negative spectral weights in
the effective theory; propagation outside the derived cone; extra modes or imposed equations in the Maxwell sector.

% --- source: tex/s27_baseline#2 (The unrestricted scalar)
\subsection{The unrestricted scalar}

\begin{proposition}[A free field on $\R^{3,3}$]
\label{prop:toy}
For the action $S_6=\tfrac12\int d^6X\,(\eta^{AB}\partial_A\phi\,\partial_B\phi-m^2\phi^2)$ in coordinates $(t,u,v,x,y,z)$, with $t,u,v$ temporal:
\begin{enumerate}[label=(\roman*),leftmargin=2.2em]
\item plane waves $e^{-i\omega t+i(q_1u+q_2v)+i\bk\cdot\mathbf x}$ obey
\begin{equation}
\label{eq:toydisp}
\omega^2=|\bk|^2+m^2-q^2,\qquad q^2=q_1^2+q_2^2 ;
\end{equation}
\item for $q^2>|\bk|^2+m^2$ the growth rate is $\gamma=\sqrt{q^2-|\bk|^2-m^2}$, and $(1+q^2)^{-s/2}e^\gamma\to\infty$ for every $s$. Plane waves are not of finite
norm on the noncompact Cauchy surface, so the argument is applied to wave packets concentrated at large surplus momentum, or to plane waves in a
periodic box: data small in any Sobolev norm produce arbitrarily large solutions at $t=1$, so the Cauchy problem with respect to $t$ is ill-posed in
the sense of Hadamard \cite{Hadamard23};
\item the canonical energy density $\tfrac12(\pi^2-\phi_u^2-\phi_v^2+|\nabla_{\mathbf x}\phi|^2+m^2\phi^2)$ averages to $\tfrac14a^2(m^2-q^2)$ on $\phi=a\sin qu$ in a
periodic box in $u$, and approximately so on wave packets concentrated near such a mode: the energy is unbounded below.
\end{enumerate}
\end{proposition}

\noindent Executable check: c2-1 (exact symbolic or integer arithmetic; \cref{app:methods}).

\begin{figure}[t]
\centering
\begin{tikzpicture}
\begin{axis}[width=10.5cm,height=6.6cm,xmin=0,xmax=3.2,ymin=-5.5,ymax=4.5,axis on top,clip=false,
  axis lines=left,
  xlabel={surplus momentum $q$},ylabel={$\omega^2$},xtick={1.4142},xticklabels={$\sqrt{|\bk|^2+m^2}$},ytick={0},
  ylabel style={rotate=-90}]
\addplot[draw=none,fill=scol!12,forget plot] coordinates {(1.4142,-5.5) (3.2,-5.5) (3.2,4.5) (1.4142,4.5)} -- cycle;
\draw[gray,dashed] (axis cs:0,0) -- (axis cs:3.2,0);
\addplot[tcol,very thick,domain=0:2.739,samples=80] {2-x^2};
\addplot[only marks,mark=*,mark size=2.6pt,gcol] coordinates {(0,2)};
\node[gcol,anchor=south west,font=\small,align=left] at (axis cs:0.06,2.2) {$q=0$: the four-dimensional\\ mass shell, kept by the constraints};
\node[scol,anchor=north west,font=\small,align=left] at (axis cs:2.25,4.2) {unstable band:\\ $\omega=\pm i\sqrt{q^2-|\bk|^2-m^2}$};
\end{axis}
\end{tikzpicture}
\caption{The dispersion relation of a free scalar on $\R^{3,3}$ at fixed spatial momentum (\cref{prop:toy}). Every surplus momentum beyond $\sqrt{|\bk|^2+m^2}$ gives
an exponentially growing mode, with a rate that grows without bound. The constraints of \cref{prop:constraints} keep exactly the point $q=0$.}
\label{fig:toy}
\end{figure}

\begin{proposition}[No time orientation in signature $(3,3)$]
\label{prop:orient}
The timelike set $\{\eta(X,X)>0\}$ of $\R^{p,q}$ is connected for $p\ge2$: the path $X(s)=(\cos s,\sin s,0;\,0,0,0)$ joins $X$ to $-X$ with $\eta=1$ throughout. Only for $p=1$
does it split into a past and a future component.
\end{proposition}

\noindent Executable check: c2-2 (exact symbolic or integer arithmetic; \cref{app:methods}).

The split metric therefore supplies no canonical future cone invariant under the identity component $SO^+(3,3)$: a cone, or a causal order, can be
introduced only by additional data. Anything called a light cone with a future has to come from such data, here a reduction to a single time. The arena already contains such a reduction, as an input. For the geodesic test particle of \cref{prop:constraints}, $H=\tfrac12(h^{\mu\nu}k_\mu k_\nu+p_u^2+p_v^2)$
with $h=\operatorname{diag}(1,-1,-1,-1)$, the mass shell $2H=m^2$ is $p_0^2=|\mathbf p|^2+m^2-p_u^2-p_v^2$, the relation \eqref{eq:toydisp} with $q=(p_u,p_v)$. The
first-class constraints $p_u=p_v=0$ remove every mode with nonzero surplus momentum, as \cref{fig:toy} shows: the unstable band
$q^2>|\mathbf p|^2+m^2$, and with it the stable modes $0<q^2\le|\mathbf p|^2+m^2$, keeping only $q=0$. They are sufficient to remove the instability
but not minimal; the minimal restriction, $q^2\le|\mathbf p|^2+m^2$, depends on the four-momentum and is not of this local form.

\begin{example}[Compact surplus times give tachyons]
\label{ex:tachyons}
If instead $u$ and $v$ are made periodic with period $2\pi R$, the surplus momenta become $n/R$ with $n\in\Z^2$, and the four-dimensional masses are
\[
m_n^2=m^2-\frac{|n|^2}{R^2}.
\]
Every mode with $|n|>mR$ is a tachyon, and there are infinitely many: for $mR=3$ there are $92$, $412$ and $1652$ of them within $|n_i|\le5$, $10$ and $20$.
Compactification, the usual cure for extra space dimensions, is no cure for extra times.
\end{example}

\noindent Executable check: c2-3, c2-4 (exact symbolic or integer arithmetic; \cref{app:methods}).

An ultrahyperbolic problem can be made well posed by a nonlocal constraint on the data \cite{CraigWeinstein09}, and the dimensionality of spacetime has been
argued from such considerations \cite{Tegmark97}. The constraints of the arena, $\partial_u\phi=\partial_v\phi=0$, are local and stronger: they keep only the zero modes,
which is dimensional reduction. Either way the cure is a restriction, and at this stage it is an input. \Cref{sec:constraints} asks where it can come from; the next
two sections ask whether the four-plane it presupposes can be selected.

% --- source: tex/s28_selector#0 ((head))


% --- source: tex/s28_selector#1 (The kinematic selector)
\section{The kinematic selector}
\label{sec:selector}

Choosing a Lorentzian four-plane inside $\R^{3,3}$ is equivalent to choosing its orthogonal complement, a two-plane, and a two-plane is described by its unit
bivector. The question is which bivectors describe a Lorentzian complement, and what the arena does with them.

% --- source: tex/s28_selector#2 (One constituent)
\subsection{One constituent}

\begin{theorem}[The selector is a complex structure]
\label{thm:selector}
Let $C=a\wedge b$ be a simple bivector of $\R^{3,3}$ and $N=\langle a,b\rangle$.
\begin{enumerate}[label=(\roman*),leftmargin=2.2em]
\item $C^2=-\det\mathrm{Gram}(a,b)\,\id$.
\item For a nondegenerate $N$ with $C$ normalised: $C^2=-\id$ exactly when $N$ is definite, that is exactly when $N^\perp$ is Lorentzian, of signature $(1,3)$
when $N$ is positive and $(3,1)$ when $N$ is negative; and $C^2=+\id$ exactly when $N^\perp$ has signature $(2,2)$. By (i), $C^2=0$ exactly when $N$ is degenerate.
\item For a nondegenerate $N$, the centraliser of $C$ in $\so{3,3}$ is $\so N\oplus\so{N^\perp}$, of dimension seven, so the choices of $N$ of a given
nondegenerate type form an eight-dimensional homogeneous space. The degenerate strata differ: for the totally null plane $\langle e_1+e_4,e_2+e_5\rangle$,
$C_0=(e_1+e_4)(e_2+e_5)$ has $C_0^2=0$ and a nine-dimensional centraliser.
\item The surplus plane is such an $N$, with $C=J=e_2e_3$ and $N^\perp$ the interior four-plane.
\end{enumerate}
\end{theorem}

\begin{proof}
(i) is the Clifford identity $(a\wedge b)^2=(a\cdot b)^2-a^2b^2$. (ii) For an orthonormal pair, $(ab)^2=-a^2b^2$, which is $-1$ when $a^2=b^2$ and $+1$ otherwise, and the
signature of $N^\perp$ is the complement of that of $N$ in $(3,3)$. (iii) holds for coordinate planes by \eqref{eq:bladecomm}, and every nondegenerate plane is an
$SO(3,3)$-image of a coordinate plane of the same signature, by Witt's theorem; the centraliser of $C_0$ is computed directly.
\end{proof}

\noindent Executable check: c1-1, c1-12 (exact symbolic or integer arithmetic; \cref{app:methods}).

For coordinate planes $C$ is moreover antisymmetric on $S$ when $C^2=-\id$, an orthogonal complex structure; for tilted planes it is not. Orthogonality refers to the
Euclidean structure of $S$, which singles out the temporal triple, while $C^2=-\id$ is invariant. The distinction will matter.

Spinors give the light cone. The volume element $\omega_S=e_4e_5e_6$ satisfies $\omega_Se_A\omega_S^{-1}=-e_A^{\!\top}$ for every generator, so
$\langle\psi,\varphi\rangle=\psi^{\!\top}\omega_S\varphi$ is a $\mathrm{Spin}(3,3)$-invariant pairing, and it pairs $S^+$ with $S^-$.

\begin{theorem}[The light cone as twisted squares]
\label{thm:cone}
Let $\psi$ range over the chiral half $S^+$.
\begin{enumerate}[label=(\roman*),leftmargin=2.2em]
\item The forms $m_A=\omega_Se_A$ are antisymmetric and satisfy $X^{\!\top}m_A+m_AX=-\omega_S[X,e_A]$ for every bivector $X$; their restrictions to $S^+$ span
$\Lambda^2(S^+)^*$, so $\psi\wedge\varphi\mapsto(\psi^{\!\top}\omega_Se_A\varphi)_A$ is an equivariant isomorphism $\Lambda^2S^+\cong\R^{3,3}$, the Klein correspondence of
\cref{thm:klein}.
\item On $S^+$, $J^2=-\id$, so $S^+\cong\C^2$, and the interior $\so{1,3}$ acts $\C$-linearly, as $\slr{2,\C}$.
\item For a unit simple bivector $C$ with plane $N$ define $v_C(\psi)_A=\psi^{\!\top}\omega_Se_AC\psi$. Then $v_C(\psi)\in N^\perp$ and $\eta(v_C,v_C)=0$ identically,
and the image has dimension three.
\item If $C^2=-\id$, the component of $v_C$ along the timelike direction of $N^\perp$ is definite: $v_C$ fills one nappe of the light cone of $N^\perp$. If
$C^2=+\id$, no component is definite, and $v_C$ fills the null vectors of a $(2,2)$ plane, which has no time orientation.
\item $v_{-C}=-v_C$; on $S^-$ the $J$-twisted square lies in the other nappe; and on $S^+$, $v_J=-j$, where $j_A=\psi^{\!\top}\gamma^0e_A\psi$ is the Dirac current.
\end{enumerate}
\end{theorem}

\noindent Executable check: c1-2, c1-3 (exact symbolic or integer arithmetic; \cref{app:methods}).

\begin{figure}[t]
\centering
\begin{tikzpicture}[>=Stealth,font=\small]
\begin{scope}
\fill[tcol!10] (0,0) -- (1.6,2.2) arc (0:180:1.6 and 0.35) -- cycle;
\draw[tcol,thick] (-1.6,2.2) -- (0,0) -- (1.6,2.2);
\draw[tcol,thick] (0,2.2) ellipse (1.6 and 0.35);
\draw[tcol!60,dashed] (-1.6,-2.2) -- (0,0) -- (1.6,-2.2);
\draw[tcol!60,dashed] (0,-2.2) ellipse (1.6 and 0.35);
\draw[->,gray] (0,-2.7) -- (0,2.9) node[above] {$\gamma^0$};
\node[tcol,align=center] at (0,3.6) {$C=J$, $C^2=-\id$: complement $(1,3)$};
\node[anchor=west,align=left] at (1.75,1.7) {$v_J(\psi)$, $\psi\in S^+$:\\ one nappe};
\node[anchor=west,align=left,tcol!70] at (1.75,-1.7) {$S^-$ or $-J$:\\ the other nappe};
\end{scope}
\begin{scope}[xshift=8.2cm]
\draw[scol,thick] (-1.6,-1.6) -- (1.6,1.6); \draw[scol,thick] (-1.6,1.6) -- (1.6,-1.6);
\draw[scol,thick,dashed] (-2.1,-0.7) -- (2.1,0.7); \draw[scol,thick,dashed] (-2.1,0.7) -- (2.1,-0.7);
\fill[scol!10] (0,0) circle (0.25);
\node[scol,align=center] at (0,3.6) {$C^2=+\id$: complement $(2,2)$};
\node[align=center] at (0,-2.6) {null vectors of a $(2,2)$ plane:\\ connected, no past or future};
\end{scope}
\end{tikzpicture}
\caption{The selector at one constituent (\cref{thm:selector,thm:cone}). A unit bivector $C$ with $C^2=-\id$ selects a Lorentzian complement, and the $C$-twisted squares of
chiral spinors fill one nappe of its light cone, the other nappe belonging to the other chirality or to $-C$. A bivector with $C^2=+\id$ selects a $(2,2)$ complement,
whose null set has no time orientation.}
\label{fig:selector}
\end{figure}

The selector of the reduction is therefore not a new field. Choosing the spacetime $N^\perp$ and choosing the complex structure $C$ are one choice, and in the arena that
complex structure is the charge $J$, the same $J$ that is the imaginary unit of the Dirac theory (\cref{thm:i}). Null directions are the squares of Weyl spinors twisted
by $J$, and the time orientation is fixed by the sign of $J$ together with the chirality. The mathematics is the Klein--Penrose correspondence, in which a null vector is
$\psi\bar\psi$ \cite{Penrose60,PenroseRindler84,WardWells}; what is specific to the arena is that $J$ already plays this role. All of this is kinematics: nothing yet says why $J$,
rather than an indefinite or degenerate bivector, is chosen. For degenerate planes $v_C$ is still null with a three-dimensional image, and for the totally null plane
$\langle e_1+e_4,e_2+e_5\rangle$ it lies in the totally null three-plane $\langle e_i+e_{i+3}\rangle$, an eigenspace of the grading root $\cK$.
A six-dimensional Dirac model of signature $(3,3)$ with a locally gauged $SO(3,3)$, broken locally to the Lorentz group of a selected time
direction, was studied by Dartora and Cabrera \cite{DartoraCabrera10}. That general architecture is prior art; what is examined here is the
identification of the selector with $J$ and the obstructions to its dynamical selection (\cref{sec:rel:times}).

% --- source: tex/s28_selector#3 (Three constituents)
\subsection{Three constituents}

At three constituents $J$ is the product of the fourteen generators of the Lorentz complement $W_3$, of signature $(8,6)$, and $\Omega$ is the volume element of $E$.

\begin{theorem}[The selector at three constituents]
\label{thm:selector3}
\begin{enumerate}[label=(\roman*),leftmargin=2.2em]
\item $J^2=\Omega^2=-\id$, $[J,\Omega]=0$ and $\Gamma=J\Omega$ with $\Gamma^2=\id$; $\Omega=-J$ on $\Gamma=+1$. Since $\Gamma^2=\id$ at every rung, any coordinate four-plane $E'$
with volume element $\Omega_{E'}$ and complementary volume element $J_{E'}$ has $J_{E'}^2=\Omega_{E'}^2=\prod_{i\in E'}e_i^2$: the volume element of the complement is a complex
structure exactly when $E'$ is Lorentzian, at every $n$.
\item The centraliser of $J$ among the $153$ bivectors has dimension $97=6+91$, so the choices form a space of dimension $56$.
\item With either $\mathrm{Spin}(9,9)$-invariant pairing, the product $\beta_{\rm odd}$ of the nine negative generators, which is antisymmetric, or the product $\beta_{\rm even}$
of the nine positive ones, which is symmetric, $\beta e_AJ$ is symmetric exactly for the fourteen directions of $W_3$: the $J$-twisted square lives on the complement,
not on $E$. The untwisted $\beta_{\rm odd}e_A$ is symmetric for all eighteen directions, so the vector sits in $\mathrm{Sym}^2S^+$, and the untwisted
square is not null at random spinors of $S^+$.
\end{enumerate}
\end{theorem}

\noindent Executable check: c1-4, c1-5 (exact symbolic or integer arithmetic; \cref{app:methods}).

The identification of the selector with $J$ survives at three constituents, but the covariant construction of the cone from spinors is a feature of dimension six. At
$n=3$ the cone of $E$ is recovered by the Dirac current $j_\mu=\psi^{\!\top}\gamma^0\gamma_\mu\psi$, which is future causal, being a sum of one-constituent currents, and
null on product states whose first factor is chiral (for the first factor $|000\rangle+|100\rangle$ it is $(2,0,0,0)$, timelike); but it uses the clock $\gamma^0$ and the transpose, that is, the positive inner product of the spinor module.

% --- source: tex/s29_typesel#0 ((head))


% --- source: tex/s29_typesel#1 (Type selection)
\section{Type selection}
\label{sec:typesel}

Kinematics identifies the selector with a complex structure. Dynamics would have to explain why the field that plays this role takes a value of the
complex-structure type at all, rather than an indefinite or degenerate one. The natural test is a potential: take the order parameter to be an unconstrained bivector field
$B\in\so{3,3}$, an adjoint scalar, ask that the potential be invariant under $SO(3,3)$, so that nothing is put in by hand, and require that its minimum set be exactly the orbit of $J$,
over an open set of couplings, with healthy fluctuations. Restriction to $S^+$ is the isomorphism $\so{3,3}\cong\slr{4,\R}$, and the invariant polynomials are the
polynomials in
\[
t_k(B)=\tr\bigl((B|_{S^+})^k\bigr),\qquad k=2,3,4,
\]
the coefficients of the characteristic polynomial; a continuous invariant function is constant on orbit closures (for the real orbits and their closures
see \cite{RichardsonSlodowy90,Jablonski12}).

\begin{proposition}[Invariants of the three types]
\label{prop:types}
Rotations of a positive plane and of a negative plane both have $(t_2,t_3,t_4)=(-4,0,4)$ and characteristic polynomial $(y^2+1)^2$; boosts of an indefinite plane have
$(4,0,4)$ and $(y-1)^2(y+1)^2$. The positive and the negative planes induce opposite orientations of $S^+=\R^4$: they are distinct $SL(4,\R)$-orbits with identical
invariants.
\end{proposition}

\noindent Executable check: c1-6 (exact symbolic or integer arithmetic; \cref{app:methods}).

So no invariant polynomial can prefer $(1,3)$ to $(3,1)$. Smooth invariant functions can separate the two orientations, since the orbits
are distinct; the stronger obstruction is the orbit closure of \cref{thm:flat}(i). The two are exchanged by the anti-isometries $\cK$ and $G$ and correspond to the two sign conventions for a
Lorentzian metric; physically this is harmless, but it means the selection is never unique. The type itself is harder.

\begin{theorem}[No clean selection below degree eight]
\label{thm:degree}
Let $V$ be an $SO(3,3)$-invariant polynomial on $\so{3,3}$, bounded below.
\begin{enumerate}[label=(\roman*),leftmargin=2.2em]
\item If $\deg V\le4$, then $V=m\,t_2+\lambda\,t_2^2+\mathrm{const}$.
\item If $\deg V\le7$, then $V$ restricted to $\{t_3=0\}$ is a polynomial of degree at most two in $t_2$ alone.
\item Hence if $J$ minimises $V$, so does $B_w=\tfrac35e_2e_3+\tfrac45e_5e_6$, which has $t_2=-4$, $t_3=0$ and $t_4=\tfrac{4804}{625}\neq4$, is not simple, and selects no four-plane.
\item In degree eight, $F=(t_2+4)^2+t_3^2+(t_4-4)^2$ vanishes exactly on the elements with characteristic polynomial $(y^2+1)^2$, and its Hessian at $J$ has rank one.
\end{enumerate}
\end{theorem}

\begin{proof}
Use test elements of $\slr{4,\R}$ with rational entries. $A_0=\operatorname{diag}(1,-1)\oplus\bigl(\begin{smallmatrix}0&-1\\1&0\end{smallmatrix}\bigr)$ and
$B_0=\bigl(\begin{smallmatrix}1&-1\\1&1\end{smallmatrix}\bigr)\oplus\bigl(\begin{smallmatrix}-1&-1\\1&-1\end{smallmatrix}\bigr)$ have $t_2=t_3=0$ and $t_4=4,-16$;
$C_0=\operatorname{diag}(-1,-3)\oplus\bigl(\begin{smallmatrix}2&-3\\3&2\end{smallmatrix}\bigr)$ has $t_2=0$ and $t_3=-120$; the complex and para-complex structures have
$(t_2,t_4)=(\mp4,4)$; $D_0=\operatorname{diag}(1,-1,2,-2)$ and two rotations of speeds $1,2$ have $(t_2,t_4)=(\pm10,34)$, all with $t_3=0$. (i) Along $sA_0$ and $sB_0$ a quartic
term $\lambda_2t_4$ would give $\lambda_2t_4s^4$ with both signs of $t_4$, so $\lambda_2=0$ for boundedness; along $sC_0$ a cubic term gives $-120\kappa s^3$, so $\kappa=0$. (ii) The top
homogeneous part must be non-negative, so an odd top degree vanishes, and in degree six the top part $e\,t_2^3+f\,t_3^2+g\,t_2t_4$ takes the values $\mp(64e+16g)$ on the
complex and para-complex structures and $\pm(1000e+340g)$ on the last two elements; non-negativity forces $g=-4e$ and then $e=g=0$, with $f\ge0$ from $C_0$. The quartic
coefficient of $t_4$ must again vanish, leaving $a\,t_2+b\,t_3+c\,t_2^2+h\,t_2t_3+f\,t_3^2$, which on $t_3=0$ is $a\,t_2+c\,t_2^2$. (iii) follows. (iv) $F\ge0$ with equality exactly
when $(t_2,t_3,t_4)=(-4,0,4)$, which fixes the characteristic polynomial by Newton's identities; at $J$ the gradients of $t_3$ vanish and those of $t_2$ and $t_4$ are
proportional, so the Hessian has rank one.
\end{proof}

\noindent Executable check: c1-7, c1-8 (exact symbolic or integer arithmetic; \cref{app:methods}).

\begin{figure}[t]
\centering
\begin{tikzpicture}
\begin{axis}[width=12cm,height=7cm,xmin=-13,xmax=8,ymin=0,ymax=10,axis lines=middle,clip=false,
  xlabel={$t_2$},ylabel={$t_4$},xtick={-4,4},ytick={4},
  xlabel style={anchor=west},ylabel style={anchor=south}]
\draw[gray,dashed] (axis cs:-4,0) -- (axis cs:-4,10);
\draw[gray,dotted] (axis cs:-4,4) -- (axis cs:0,4);
\draw[gray,dotted] (axis cs:-4,7.6864) -- (axis cs:0,7.6864);
\addplot[only marks,mark=*,mark size=3pt,tcol] coordinates {(-4,4)};
\addplot[only marks,mark=*,mark size=3pt,scol] coordinates {(4,4)};
\addplot[only marks,mark=square*,mark size=2.6pt,ocol] coordinates {(-4,7.6864)};
\node[tcol,anchor=east,font=\small,align=right] at (axis cs:-4.3,4) {$J$: complex structure,\\ Lorentzian complement};
\node[scol,anchor=north,font=\small,align=center] at (axis cs:4,3.6) {$\cK$-type boost:\\ $(2,2)$ complement};
\node[ocol,anchor=east,font=\small,align=right] at (axis cs:-4.3,7.69) {$B_w=\tfrac35e_2e_3+\tfrac45e_5e_6$, $t_4=\tfrac{4804}{625}$:\\ not simple, no four-plane};
\node[gray,anchor=east,font=\small] at (axis cs:-4.3,1.2) {degree $\le7$: blind along $t_2=-4$};
\end{axis}
\end{tikzpicture}
\caption{Type selection in the plane of the invariants $(t_2,t_4)$ at $t_3=0$ (\cref{thm:degree}). A bounded-below invariant potential of degree at most seven depends on $t_2$
alone along $t_3=0$, so it cannot distinguish the complex structure $J$ from the non-simple element $B_w$ on the same vertical line. Degree eight is the first that can.}
\label{fig:typesel}
\end{figure}

Even where a potential does isolate the type, the minimum cannot be clean, and the reason holds at every degree.

\begin{theorem}[Degeneracy and flat directions at every degree]
\label{thm:flat}
Let $V$ be any $SO(3,3)$-invariant function on $\so{3,3}$.
\begin{enumerate}[label=(\roman*),leftmargin=2.2em]
\item Let $n=(e_1+e_4)e_5$, a null rotation of the interior four-plane. Then $[J,n]=0$ and $n^2=0$, so $J+n$ has the invariants of $J$, but its centraliser has dimension three
and does not contain $\so{1,3}$: it breaks Lorentz invariance. With $K=\tfrac12e_1e_4$, $[K,n]=-n$ and $[K,J]=0$, so conjugation by $e^{tK}$ maps $J+n$ to $J+e^{-t}n$. Hence $J$
lies in the closure of the orbit of $J+n$, and $V(J+n)=V(J)$ whenever $V$ is continuous.
\item If $V$ is smooth and has a local minimum at $J$, its Hessian vanishes on the six $\so{1,3}$ directions. The $\Stab(J)$-invariant symmetric forms on $\so{1,3}$ are the real
and imaginary parts of the complex Killing form of $\slr{2,\C}$, and each nonzero one has signature $(3,3)$, so a positive semidefinite one is zero.
\end{enumerate}
At best the fluctuations about $J$ are one massive radial mode, eight Goldstone modes and six massless modes that are not Goldstone modes:
non-Goldstone flat directions, whose vanishing masses are enforced by invariance and stability (ii), not by a broken symmetry.
\end{theorem}

\noindent Executable check: c1-9 (exact symbolic or integer arithmetic; \cref{app:methods}).

\begin{proposition}[Kinetic terms]
\label{prop:kineticB}
Take $B$ an unconstrained adjoint scalar with a factorised kinetic term, a spacetime contraction of $\partial B$ with $\partial B$ times an invariant
symmetric form on $\so{3,3}$. Since $\so{3,3}$ is simple, that form is a multiple of the Killing form, which on the fifteen bivectors has signature $(9,6)$, boosts with
one sign and rotations with the other. At $J$ the radial direction is a rotation, the eight Goldstone directions are two rotations and six boosts, and the six $\so{1,3}$
directions are three of each. Whatever the overall sign, at least six directions are ghosts, and the induced metric on the eight-dimensional space of choices has
signature $(6,2)$.
\end{proposition}

\noindent Executable check: c1-10 (exact symbolic or integer arithmetic; \cref{app:methods}).

In covariant terms the Goldstone directions form a charged four-vector and the $\so{1,3}$ directions a two-form, and healthy vectors and two-forms owe their positivity to a
gauge symmetry that removes the wrong-sign components; a scalar order parameter has none. This is \cref{thm:ghost} again, in a new setting. There is a positive
alternative, and it shows what is missing.

\begin{proposition}[The positive alternative]
\label{prop:positivealt}
On the fifteen bivectors the positive form $\langle X,Y\rangle=\tr(XY^{\!\top})/8$ is the identity, and its invariance algebra is the six compact rotations $\so3_T\oplus\so3_S$.
Split $B=B_T+B_S+B_p$ into these two parts and the nine boosts. The potential
\[
V_+=\bigl(|B_T|^2-1\bigr)^2+|B_S|^2+|B_p|^2
\]
has as minimum set the unit sphere of $\so3_T$. Every element there is a rotation of a plane $N$ in the temporal triple, the clock is the axis of the rotation, the spacetime is the
clock line plus the spatial triple, and $J=\gamma^0\omega_T$. The Hessian at $J$ has eigenvalues $8$, $0$, $0$ and $2$ (twelve times): two Goldstone modes and everything else massive
with positive kinetic energy. But the boosts $e_1e_4,e_1e_5,e_1e_6$ of the selected spacetime are not symmetries of this sector.
\end{proposition}

\noindent Executable check: c1-11 (exact symbolic or integer arithmetic; \cref{app:methods}).

The requirement was a clean selection of the Lorentzian type over an open set of couplings, with every mode healthy and accounted for. An unconstrained $SO(3,3)$-invariant adjoint scalar
order parameter cannot deliver it: below degree eight it does not isolate the type; at every degree the minimum is degenerate with Lorentz-violating configurations and has
six non-Goldstone flat directions; and, with a factorised kinetic term, it has ghosts. The positive route selects cleanly only by giving up the boosts it was meant to produce.

% --- source: tex/s30_constraints#0 ((head))


% --- source: tex/s30_constraints#1 (Where the surplus constraints come from)
\section{Where the surplus constraints come from}
\label{sec:constraints}

The baseline showed that the constraints $p_u=p_v=0$ remove every nonzero surplus momentum, the whole unstable band together with the stable
surplus modes, and the curved model showed that they are first class. A first-class constraint
generates a gauge symmetry, so the question is which symmetry, and whether it is a natural one.

\begin{theorem}[A worldline gauge symmetry]
\label{thm:wlgauge}
For arbitrary $A_\mu(x)$ and flat $h$, with $C=A_\mu h^{\mu\nu}k_\nu$, the constraints satisfy $\{p_u,H\}=-Cp_v$, $\{p_v,H\}=Cp_u$, $\{p_u,p_v\}=0$ and $\{\ell,H\}=0$. The extended action
\[
\int d\tau\,\bigl(p\cdot\dot x-N(H-\tfrac12m^2)-\lambda_1p_u-\lambda_2p_v\bigr)
\]
is invariant under
\[
\delta u=\varepsilon_1(\tau),\qquad \delta v=\varepsilon_2(\tau),\qquad \delta\lambda_1=\dot\varepsilon_1+NC\varepsilon_2,\qquad \delta\lambda_2=\dot\varepsilon_2-NC\varepsilon_1,
\]
all other variables fixed: local translations along the normal plane, which are gauge symmetries although they are not isometries when $F\neq0$, the algebra closing with
structure functions \cite{HenneauxTeitelboim92}. Constraint preservation is thereby shown for a test worldline in an arbitrary background $A_\mu$; for interacting fields, with the
normal-plane modes dynamical, it is not established here. In Dirac quantisation the physical states are the functions $\Psi(x)$ on spacetime, and on them $\hat H\Psi=-\tfrac12\Box_h\Psi$;
the potential $A$ drops out.
\end{theorem}

\noindent Executable check: c2-5, c2-6, c2-7 (exact symbolic or integer arithmetic; \cref{app:methods}).

First-class constraints that remove surplus time directions by generating a worldline gauge symmetry are also the mechanism of two-time physics,
where a local $Sp(2,\R)$ acts on position and momentum in signature $(d,2)$ \cite{BarsDelidumanAndreev98}; the comparison, and where it ends, is
made in \cref{sec:rel:times}.

\begin{proposition}[Gauging the isometry is not enough]
\label{prop:isometryonly}
In polar normal coordinates $\ell=-p_\phi$ and $p_u^2+p_v^2=p_r^2+p_\phi^2/r^2$. Gauging only the rotation $\ell$, the isometry of the normal plane, leaves
$p_0^2=|\mathbf p|^2+m^2-p_r^2$: the radial surplus momentum keeps the unstable band. For fields, the rotation-invariant modes $J_0(qr)e^{-i\omega t+i\bk\cdot\mathbf x}$ obey
$\omega^2=|\bk|^2+m^2-q^2$.
\end{proposition}

\noindent Executable check: c2-9 (exact symbolic or integer arithmetic; \cref{app:methods}).

The symmetry that removes the instability is therefore the full group of local translations along the normal plane, not its rotation subgroup.

\begin{corollary}[Weak emergence for scalar matter]
\label{cor:weak}
If the normal translations are gauged, the scalar sector reduces to the Klein--Gordon field on $(M^4,h)$; the physical states do not depend on
$u,v$, and they are not normalisable functions on $\R^{3,3}$, on which constants in $u,v$ have infinite norm. For flat $h$ (the setting of
\cref{prop:isometryonly}), a chosen time orientation and $m^2\ge0$, the one-particle space is that of the positive-frequency solutions with the
Klein--Gordon inner product, and on its bosonic Fock space the energy is bounded below, the evolution is unitary, and $[\phi(x),\phi(y)]=0$ at
spacelike separation for the cone of $h$ \cite{WeinbergQFT}. For a globally hyperbolic curved $h$ the algebraic Klein--Gordon construction gives
microcausality from the support of the advanced-minus-retarded propagator; a preferred positive-frequency ground state and a unitary
implementation in a chosen Fock representation need further assumptions \cite{FewsterVerch12}. In the free flat product the zero modes $\psi$ and the surplus modes $\chi$ do not
mix, since translation invariance along the normal plane makes the mixed kernel $K_{\psi\chi}$ vanish, so the Schur complement is trivial and $K_{\rm eff}=K_{\psi\psi}$;
the rotating normal metric with a general $A$ is not translation invariant, and there the reduced operator itself is needed; the content is not in
the kernel but in whether $\chi$ is physical, and under the gauge principle it is not. The causal cone is that of $h$, the metric of the selected four-plane: it is inherited
from the reduction, not derived from the six-dimensional sum.
\end{corollary}

What the reduction leaves over is the geometry of the normal plane itself, and it is not healthy.

\begin{theorem}[Every normal-plane vector is a ghost]
\label{thm:ghostvec}
For $ds^2=h+\sigma\bigl[(du+B^u)^2+(dv+B^v)^2\bigr]$ with flat $h$, $B^u=b_1x^1dx^2$ and $B^v=b_2x^0dx^3$, one finds $R_6=-\tfrac\sigma4\sum F_{\mu\nu}F^{\mu\nu}$, and for the
rotation ansatz $B^u=vA$, $B^v=-uA$, $R_6=-\sigma\,r^2F^2/4$, which at $\sigma=+1$ is the ambient scalar curvature of \cref{thm:ein6}. The coefficient is odd in $\sigma$, so the
healthy sign is that of a normal plane of the sign of space, $\sigma=-1$, the ordinary Kaluza--Klein case. By \cref{thm:selector}(ii) the normal plane of a Lorentzian
four-plane in $\R^{3,3}$ always has the sign of time. Relative to the ambient Einstein--Hilbert action and these ans\"atze, every Kaluza--Klein vector from the normal plane, from its
translations or its rotation, therefore gravitates with the wrong sign, and by \cref{thm:wlgauge} none of them couples to spinless matter in the
zero-normal-momentum sector. The theorem concerns normal vectors obtained by this Einstein--Hilbert and Kaluza--Klein route.
\end{theorem}

\noindent Executable check: g2-7 (exact symbolic or integer arithmetic; \cref{app:methods}).

The failure condition ``ghosts or negative spectral weights'' is therefore met by the normal-plane vectors, and the photon cannot come from them. If the Maxwell field
emerges, it must come from elsewhere: from the univectorial potential of the interior reading, or from the field-strength reading $\R^{3,3}=\Lambda^2\R^{1,3}$, where the
energy is the positive Hermitian form $|\bB|^2+|\bE|^2$ (\cref{sec:spin}).

\section{The clock reduction and the parameter \texorpdfstring{$\tau$}{tau}}
\label{sec:tau}

Every statement of this manuscript about a four-momentum, an observable four-plane or the absence of uphill decays rests on
\cref{assum:clockred}: physical states and observables are those of the reduced shell of one clock $\hat n$. This section says what the
assumption does, how far its implementation reaches, and what the parameters along which things evolve are.

\subsection{What the reduction is, and how far its implementation reaches}

The assumption is implemented by the first-class constraints $p_u=p_v=0$ of \cref{prop:constraints}, gauged by the local normal translations of
\cref{thm:wlgauge}; for fields it is the restriction $\partial_u\phi=\partial_v\phi=0$, which is dimensional reduction (\cref{sec:baseline}). Four
qualifications belong to it.
\begin{enumerate}[label=(\roman*),leftmargin=2.2em]
\item It is \emph{imposed}. The gauge principle is consistent but chosen: the multipliers $\lambda_1,\lambda_2$ are added to the worldline action by hand
(\cref{thm:wlgauge}), and \cref{sec:typesel} shows that no invariant potential for an unconstrained adjoint bivector order parameter has
the orbit of $J$ as its minimum set, and that the factorised invariant kinetic term has ghosts; other representations, constrained order
parameters, gauge or soldering constructions and derivative couplings are not excluded.
\item It is \emph{sufficient, not minimal}. The constraints remove every nonzero surplus momentum, the stable modes $0<q^2\le|\mathbf p|^2+m^2$
together with the unstable band (\cref{fig:toy}); the minimal restriction depends on the four-momentum and is not of this local form.
\item Its \emph{preservation} is shown for a test worldline in an arbitrary background $A_\mu$ and a flat $h$ (\cref{thm:wlgauge}). For
interacting fields with dynamical normal-plane modes it is not established.
\item It removes the surplus times and leaves the normal-plane \emph{geometry}: relative to the ambient Einstein--Hilbert action, every Kaluza--Klein
vector of the normal plane is a ghost (\cref{thm:ghostvec}), and none couples to spinless matter in the zero-normal-momentum sector.
\end{enumerate}

\subsection{The worldline parameter}

\begin{proposition}[The worldline parameter is a gauge parameter]
\label{prop:wltau}
For the extended action $S=\int d\tau\,(p\cdot\dot x-N(H-\tfrac12m^2)-\lambda_1p_u-\lambda_2p_v)$ of \cref{thm:wlgauge}:
\begin{enumerate}[label=(\roman*),leftmargin=2.2em]
\item $S$ is invariant under $\tau\mapsto\tau'=f(\tau)$ with $f'>0$, $x$ and $p$ transforming as scalars and $N,\lambda_1,\lambda_2$ as densities,
$N'=N/f'$, $\lambda_i'=\lambda_i/f'$. Varying $N$ gives the mass shell $2H=m^2$, a first-class constraint; on shell the reparametrisations are
generated by the total first-class Hamiltonian $N(H-\tfrac12m^2)+\lambda_1p_u+\lambda_2p_v$, multipliers included.
\item On the reduced shell $p_u=p_v=0$, $h^{\mu\nu}p_\mu p_\nu=m^2>0$, the proper time of $h$ along the worldline is $d\tau_h=|N|m\,d\tau$.
\end{enumerate}
Hence $\tau$ is not an observable. For $m>0$, on a segment where the lapse does not vanish, the gauge $N=1/m$ (if $N>0$) or $N=-1/m$
(if $N<0$; reparametrisations with $f'>0$ preserve the sign of $N$) makes it the proper time of $h$; a lapse crossing zero needs separate treatment, and the time read by the clock is the proper time of the
selected four-plane, not a seventh coordinate.
\end{proposition}

\begin{proof}
(i) Under the reparametrisation $p\cdot\dot x\,d\tau$ is invariant, and $N(H-\tfrac12m^2)\,d\tau=N'(H-\tfrac12m^2)\,d\tau'$, $\lambda_ip_i\,d\tau=\lambda_i'p_i\,d\tau'$.
The variation of $N$ gives $H-\tfrac12m^2=0$; with $-NH$ alone it would give $H=0$, the massless shell. The constraint
algebra of \cref{thm:wlgauge} contains $\{p_u,H\}=-Cp_v$, $\{p_v,H\}=Cp_u$ and $\{p_u,p_v\}=0$, so $H-\tfrac12m^2$, $p_u$, $p_v$ are first class. (ii) With
$k_\mu=p_\mu-A_\mu\ell$ and $\ell=vp_u-up_v$, Hamilton's equation for $x^\mu$ is $\dot x^\mu=Nh^{\mu\nu}k_\nu$. On the reduced shell $\ell=0$, so $k=p$ and
$h_{\mu\nu}\dot x^\mu\dot x^\nu=N^2h^{\mu\nu}p_\mu p_\nu=N^2m^2$, so $d\tau_h=|N|m\,d\tau$.
\end{proof}

\noindent Executable check: c2-8 (exact symbolic or integer arithmetic; \cref{app:methods}).

This is the standard reparametrisation invariance of a relativistic worldline \cite{HenneauxTeitelboim92}, written for the reduced arena. It
says what $\tau$ is and what it is not: a label of points on a worldline, fixed only by a gauge choice, whose physical content is the proper
time of $h$. It does not turn the six-dimensional manifold, or its tangent fibres, into a seven-dimensional spacetime.

\subsection{The interior parameter}

The interior dynamics of \cref{ch:ch04_rigidbody}, the shifted operators $\partial_\tau\pm\kappa G$, the rotor and the closure, also evolves in a
parameter written $\tau$, with the dimension of a time. That parameter is not the worldline parameter of \cref{prop:wltau}, and this
manuscript does not relate it to the time read by a clock (\cref{sec:esmap:map}). Two further steps would be needed to do so. The rate
$\kappa=c/\ell$ of the binding would have to be identified with a rest frequency, $mc^2=\hbar\kappa$, which is an added hypothesis
(\cref{sec:restscale}); and the rotor would have to be quantised, with its angle related to an external time. Both are open. Where the
interior parameter is read as Newtonian time, in the orbital extension of \cref{sec:orbital}, the reading is part of the identification and is
stated there.

\subsection{Local clock rotations}
\label{sec:localclock}

All unit clocks are related by rotations of $V^+$, and the rotor carrying $e_1$ to $\widehat W$ conjugates the free Dirac data (\cref{rem:clockcov}).
That is a global covariance of free data. A position-dependent rotor $R(x)$ adds the term $R^{-1}\partial_\mu R\in\so3_+$ to the Dirac operator, so a
local clock symmetry needs a connection with values in the rotations of $V^+$ and its transformation law. Since a rotation of the temporal triple
alone changes the binding root $G$ (\cref{rem:clockcov}(b)), such a connection would also have to transport $G$, or be restricted to the matched
rotations of both triples. It is not constructed here. Nor is it the gauge symmetry of \cref{thm:wlgauge}: that symmetry translates
along the normal plane of a fixed four-plane and does not rotate the clock. Global free equivalence therefore establishes neither a local clock
gauge symmetry, nor a dynamical selection of the clock, nor a well-posed interacting theory with three times.

\subsection{What survives the reduction}

Under the reduction every momentum lies in $E_{\hat n}=\R\hat n\oplus V^-$, of signature $(1,3)$, and the following hold.
\begin{enumerate}[label=(\roman*),leftmargin=2.2em]
\item \emph{Kinematics.} The six-dimensional mass equals the mass read by the clock, $M_6=m_4(\hat n)$, and the ordinary four-dimensional
decay inequality holds: no uphill decay survives (\cref{e7:prop:uphill}(iv)). The uphill decays and the opened channels of
\cref{e7:prop:uphill}(ii)--(iii) use temporal momenta transverse to the clock, which the reduction removes; they are warnings about the unreduced
kinematics, not predictions, and the ambient temporal momenta are not observable energies (\cref{ch:ch10_classical}).
\item \emph{Scalar matter.} With the normal translations gauged, the scalar sector is the Klein--Gordon field on $(M^4,h)$, with energy bounded
below, unitary evolution and microcausality for the cone of $h$ for flat $h$; on a globally hyperbolic curved $h$, microcausality of the algebraic
construction, with a preferred state and unitarity needing further assumptions (\cref{cor:weak}).
\item \emph{Removed.} Every surplus-momentum mode, stable or not; with compact surplus times the tower of tachyons of \cref{ex:tachyons} is
removed with them.
\item \emph{Not removed.} The ghost vectors of the normal plane (\cref{thm:ghostvec}) and the type-selection obstructions of
\cref{sec:typesel}, which concern the four-plane the reduction presupposes.
\end{enumerate}

\section{A universal evolution parameter: an open hypothesis}
\label{sec:universaltau}

A universal invariant evolution parameter is not new. In the parametrised relativistic dynamics of Horwitz and Piron \cite{HorwitzPiron73},
an invariant parameter evolves events in a fixed spacetime, and for a many-body system it is taken to be universal, as Newtonian time is
\cite{HorwitzPiron73,Horwitz20}. What may be specific to this framework is the connection of
such a parameter with the clock reduction of \cref{sec:tau}: the parameter would act on the split six-dimensional arena, and the time read
by a clock would follow from it through the reduction to the selected four-plane. The hypothesis is stated in a limited form, so that its
requirements are explicit. It is not a result.

\begin{quote}
\emph{Universal-evolution hypothesis.} There is one parameter $\tau_U$, external to $\R^{3,3}$, along which all fields of the arena evolve, and
the time read by a clock is obtained from $\tau_U$ through the clock reduction.
\end{quote}

As in the theory of Horwitz and Piron, the parameter is not a coordinate of spacetime. An external
parameter does not turn $\R^{3,3}$, or its tangent fibres, into a seven-dimensional physical spacetime. A proposal would have to specify the
following.
\begin{enumerate}[label=(\arabic*),leftmargin=2.2em]
\item \emph{What evolves}: fields on $\R^{3,3}$, fields on the reduced four-plane, the interior variables of \cref{ch:ch04_rigidbody}, or all of
them, and in which configuration space.
\item \emph{The action and the constraint equations}, including the constraints that implement the clock reduction, and their preservation under
the $\tau_U$-evolution with interactions.
\item \emph{Reparametrisation properties.} If the action is invariant under $\tau_U\mapsto f(\tau_U)$, its generator is a constraint and $\tau_U$ is
not observable, as for the worldline parameter of \cref{prop:wltau}. If it is not, $\tau_U$ need not be a preferred time that breaks the Lorentz
algebra of the selected four-plane or $SO(3,3)$: the action $\int d\tau_U\,(p\cdot dX/d\tau_U-\eta^{AB}p_Ap_B/2M)$ is $SO(3,3)$-invariant with $\tau_U$
a scalar, as in the theory of Horwitz and Piron. Which symmetries are kept, and which broken, has to be stated.
\item \emph{The relation to measured time}: how the proper time of $h$, which is what a clock reads (\cref{prop:wltau}(ii)), follows from $\tau_U$.
\item \emph{Causal observables}: microcausality with respect to a stated cone. A self-adjoint generator gives unitary evolution in $\tau_U$
even when it is unbounded below, so unitarity and the existence of a ground state, which needs a generator bounded below, are separate
requirements.
\item \emph{Compatibility with the obstructions of this chapter.} The type-selection results of \cref{sec:typesel} concern potentials and
$SO(3,3)$-invariant kinetic terms, and they hold whatever the evolution parameter. The ghost vectors of \cref{thm:ghostvec} come from the
six-dimensional curvature. The ill-posedness of \cref{prop:toy} is a statement about evolution in a temporal coordinate. A first-order evolution
in $\tau_U$ can pose a different Cauchy problem, but its generator then has to be bounded below for a ground state to exist, and an
$SO(3,3)$-invariant second-order operator such as $\Box_\eta$ is not, since its symbol is indefinite.
\end{enumerate}
Until a mechanism meeting these requirements is exhibited, the hypothesis remains open.

\section{What the causal analysis establishes}
\label{sec:causalverdict}

\Cref{tab:causal} collects the outcome against the failure conditions fixed at the start of \cref{sec:baseline}, together with the status of the clock reduction and of the evolution parameters.

\begin{table}[t]
\centering
\small
\begin{tabular}{@{}>{\raggedright\arraybackslash}p{2.3cm}>{\raggedright\arraybackslash}p{6.6cm}>{\raggedright\arraybackslash}p{5.1cm}@{}}
\toprule
stage & result & verdict\\
\midrule
baseline & growing modes and no ground state; the surplus constraints remove the band together with every stable surplus mode; compactification makes tachyons & obstruction confirmed, sufficient cure located\\
reduction & kinematics: selector $=J$, cone $=$ twisted squares at one constituent; dynamics: no mechanism from an unconstrained $SO(3,3)$-invariant adjoint scalar; the positive route breaks the boosts &
kinematics exact; dynamical selection fails as posed\\
effective theory & scalar matter: no mixing in the flat product, $K_{\rm eff}=K_{\psi\psi}$; normal-plane vectors are ghosts (Einstein--Hilbert route) & failure condition met by the vectors\\
causality & the scalar sector is microcausal for the cone of $h$ (flat $h$; globally hyperbolic $h$ for the algebraic construction) & conditional on gauged translations; inherited\\
Maxwell & normal-plane vectors excluded as photons & blocked on this route\\
clock reduction & \cref{assum:clockred} implemented by first-class constraints for a test worldline; removes stable surplus modes too &
imposed; sufficient, not minimal; interacting case open\\
parameters & the worldline $\tau$ is a gauge parameter whose content is the proper time of $h$; the interior $\tau$ is not related to a clock &
fixed (worldline); open (interior)\\
uphill decays & none after the reduction; the unreduced examples are warnings & conditional on the reduction\\
local clock symmetry & needs a connection in $\so3_+$ that also transports $G$ & not constructed\\
universal $\tau_U$ & prior art (Horwitz--Piron); no mechanism in this arena; requirements listed & open hypothesis\\
\bottomrule
\end{tabular}
\caption{The causal analysis against the failure conditions fixed in advance (\cref{sec:baseline,sec:selector,sec:typesel,sec:constraints,sec:tau,sec:universaltau}).}
\label{tab:causal}
\end{table}

Weak emergence holds for scalar matter, conditionally on a gauge principle, local translations along the normal plane, that is consistent but chosen rather than derived.
Strong emergence through an unconstrained $SO(3,3)$-invariant adjoint scalar order parameter is excluded; whether the quantum dynamics of the
whole system can select the Lorentzian sector, the motivation stated at the opening of the chapter, remains open. What survives is a sharper statement of the question. The selector is the complex structure $J$, and every
construction here that produced a healthy Lorentzian sector also used one extra datum: a positive inner product, which is the transpose of the realisation, equivalently
the temporal triple. Within these computations the totality of $SO(3,3)$-invariant dynamics fixes at most the type of the structure, never the instance, and it fixes even
the type only imperfectly. The datum that singles out time, wherever things worked, is the positivity of an inner product; in quantum mechanics it is the positivity of
the Hilbert-space norm. That this positivity fixes causal order is a conjecture; it indicates where a derivation of causal order from the split
geometry would have to look.

Several routes remain open. A gauge mechanism for type selection, an $SO(3,3)$ analogue of the MacDowell--Mansouri and Stelle--West breaking
\cite{MacDowellMansouri77,StelleWest80}, $SO(3,3)\to SO(2)\times SO(1,3)$, might make the would-be ghosts gauge and solder the eight broken directions to the base; the six flat directions are not gauge directions at
linear order, and removing them would need further constraints or a special action. A principle that makes the normal translations gauge, rather than adding the multipliers by hand, would turn weak emergence from a conditional statement into a
derived one. At three constituents, a $\mathrm{Spin}(9,9)$-covariant construction of the cone of $E$, for which pure spinors are the natural candidates, or a proof that
none exists without the positive inner product, would settle whether the cone survives the tower.

